스키마에 부합하는 게임 내 행동 제안이라도 접근할 수 없는 물체를 요구하거나 권한이 없는 상태 전이, 혹은 현재 시점에서 공개할 수 없는 비밀 발화를 포함할 수 있다. 서사적 정합성(grounding)과 개연성 있는 대화 생성 \cite{S03_weir2024ontologically,S20_ashby2023personalized}은 후보 콘텐츠를 제공할 뿐이며, 게임 엔진은 현재 상태에서 어떤 전이를 승인할지 별도로 판정해야 한다.

\system{}는 이 판정을 외부의 제안--커밋(proposal-to-commit) 경계에서 강제한다(그림~\ref{fig:architecture}). 모델과 어댑터가 구조화된 후보를 제출하면 결정론적 게이트가 행동 정책, 선행조건, 물체 도달 가능성, NPC 지식, 정보 공개 권한, 효과, 퀘스트 단계 단조성을 검사한다. 검사에 실패하면 게이트는 제한된 수리를 위한 구조화된 오류를 반환하며, 수리는 정본 상태(canonical state)가 아닌 후보만 수정하고 동일 상태 재검증을 통과한 경우에만 상태를 갱신한다. 여덟 역할 읽기 전용 온톨로지 뷰는 쓰기 권한 없이 출처 선언과 검사·결정을 연결하며, 해시 연결 전이 기록은 의미적 리플레이와 배정 전수 회계(assignment-complete accounting)를 보장한다.

개념적으로 이 계약은 단순 사실 일치(factual grounding)를 넘어 NPC의 \emph{인식론적 지식 경계}(행위자가 알고 있는 사실)와 \emph{정보 공개 권한 정책}(현재 퀘스트 단계에서 발화할 수 있는 사실)을 분리한다. 방법론적으로는 LLM 호출이 0회인 결정론적 검증 계층(게이트 검사, 유도 수리, 엔진 리플레이)과 실패를 누락 없이 보존하는 제한된 확률적 제안기 계층을 엄격히 분리한다.

결정론적 감사에서 스키마 전용 대조군은 32개 유한 그래프 조합에 걸쳐 기준상 무효인 제안 260/416건을 커밋한 반면, 모든 의미 게이트 조건은 0/416건을 기록했고 최초 유효 제안 156/156건(오탐 0/156) 및 고유 구조 서명 11/11종을 모두 보존했다. 반례 유도 수리는 최초 무효 제안 260건 중 128건을 유효 커밋으로 복구했다. 단일 홉 정보 공개 검사에서도 일반 규칙과 온톨로지 게이트는 60/60건 일치하며 위험 공개 51건을 차단했고, 후속 48건 다중 홉 폐쇄 벤치마크에서는 양자를 뚜렷이 구분한다(\texttt{ontology\_}\allowbreak\texttt{multihop\_}\allowbreak\texttt{gate} 48/48 대 \texttt{flat\_}\allowbreak\texttt{1hop\_}\allowbreak\texttt{rule} 20/48, 위험 누출 28건 차단). 여기에 확장 0-LLM 감사, 불변식 \textbf{I1--I4}에 대한 형식 정리 1--4, 24에피소드 라이브 다중 퀘스트 Gemini 평가(2종 모델, 3개 월드, 위험 커밋 0건, 엄격 리플레이 24/24), 2차 한·영 우회 검증기(우회 재현율 97.2\% 및 양성 오거부 0.0\%), 그리고 실제 트레이스 결속 36항목 블라인드 인간 평가 패킷(\texttt{AWAITING\_}\allowbreak\texttt{HUMAN\_}\allowbreak\texttt{RATERS})과 Croissant~1.0 메타데이터를 더해 재현 가능한 공개 경계를 완성한다.

\section{관련 연구와 연구 공백}
정합성, 기억, 계획, 역할 기반 에이전트 \cite{S03_weir2024ontologically,S20_ashby2023personalized,S07_park2023generative,S08_shao2023character} 및 플레이어·대화 시스템 \cite{S21_akoury2023player,S22_hochreiter2026beyond,S23_yin2026contextualized,S02_figueiredo2025scaffolded}은 맥락적 상호작용을 지원하는 반면, 본 연구는 인코딩된 현재 상태와 공개 정책이 제안된 전이를 허용하는지를 판정한다. 정서 및 몰입 모델링 \cite{S24_yannakakis2011experiencedriven,S25_melhart2020engagement,S26_wang2026engagement}은 생성 스타일을 안내할 뿐 상태 변경 권한을 부여하지 않는다.

기존 벤치마크와 자동 플레이테스팅은 상태, 체화 환경, 콘텐츠, 과정, 사회적 상호작용, 장기 과업을 다루며 \cite{S11_cote2019textworld,S12_wang2022scienceworld,S13_fan2022minedojo,S14_perezliebana2019gvgai,S15_liu2024agentbench,S16_ma2024agentboard,S17_zhou2024sotopia,S18_paglieri2025balrog,S19_lepelletier2022playtesting}, NCP-Bench는 서사적 일관성과 사실 정합성을 평가한다 \cite{S58_ma2026ncpbench}. 문법·스키마·어포던스·규칙 기반 기법 \cite{S09_geng2023grammar,S10_beurerkellner2023prompting,S36_scholak2021picard,S50_kambhampati2024llmmodulo,S51_ichter2023saycan,S52_wang2025agentspec,S53_li2025settingdc,S54_zhou2025story2game,S42_buongiorno2024pangea}, 반례 기반 수리 \cite{S45_madaan2023selfrefine,S48_shinn2023reflexion,S46_chen2023selfdebug,S47_solarlezama2006sketching}, 재현성·리플레이 표준 \cite{S31_agarwal2021precipice,S33_pineau2021reproducibility,S56_feng2025agentrr,S59_sun2025bdc}이 본 설계의 토대다. 스키마 적합성과 행위 합법성, 또는 사실 정합성과 발화 안전성을 동일시하던 선행 연구와 달리, \system{}는 스키마 파싱, 인식론적 지식, 공개 권한, 정본 상태 커밋을 독립 감사 가능한 게이트로 분리한다.
